% Generated draft: requires mathematical and source review.
\documentclass[11pt]{article}
\usepackage[margin=25mm]{geometry}
\usepackage{fontspec}
\setmainfont{Latin Modern Roman}
\usepackage{amsmath,amssymb,mathtools,unicode-math}
\setmathfont{Latin Modern Math}
\usepackage{enumitem,needspace,xurl,hyperref}
\setlength{\parindent}{0pt}
\setlength{\parskip}{5pt}
\setlength{\emergencystretch}{3em}
\newcommand{\N}{\mathbb{N}}
\newcommand{\Z}{\mathbb{Z}}
\newcommand{\Q}{\mathbb{Q}}
\newcommand{\R}{\mathbb{R}}
\newcommand{\C}{\mathbb{C}}
\newcommand{\F}{\mathbb{F}}
\newcommand{\A}{\mathbb{A}}
\newcommand{\PP}{\mathbb{P}}
\newcommand{\E}{\mathbb{E}}
\newcommand{\eps}{\varepsilon}
\newcommand{\dd}{\,\mathrm{d}}
\begin{document}
\section*{UnsolvedMath ID: 701220}
\textbf{Source problem code:} catalogue entry
\subsection{Definitions and mathematical statement}
Let \(F\) be a field of characteristic \(p>0\), and let \(G\) be a connected linear reductive algebraic group over \(F\). The source does not specify whether its notation \(G\) denotes \(F\)-rational points or geometric points over an algebraic closure. In the formulas below, \(G\) denotes the group of points intended by the source; this choice remains unresolved. Closed and open subsets are understood in the Zariski topology on that point set.

For a subset \(X\subseteq G\) and an integer \(k\geq 1\), define
\[
X^k=\{x_1\cdots x_k\mid x_1,\ldots,x_k\in X\}.
\]
The source asks the following two questions, with the field and algebraic group universally quantified as above.
\begin{enumerate}
\item[(a)] For every closed subset \(X\subseteq G\), does there exist an integer \(c=c(X)>1\) such that \(X^c\) is closed?
\item[(b)] For every conjugacy class
\[
X=\{gxg^{-1}\mid g\in G\},\qquad x\in G,
\]
if \(X^2\) contains an open subset of \(G\), must \(X^2=G\)?
\end{enumerate}
The source does not say that the open subset in (b) is nonempty. Literally, the hypothesis allows the empty open subset and is therefore automatic. Requiring a nonempty open subset would give a different, potentially intended formulation; that requirement has not been silently inserted here. The point domain also affects the meaning of a conjugacy class and must be clarified before a definitive formulation is accepted.
\subsection{Short English statement}
For a connected reductive linear algebraic group in positive characteristic, (a) asks whether every closed subset has a closed product power with exponent greater than one. Part (b) asks whether the square of a conjugacy class containing an open subset equals the whole group. The intended point domain and the possible requirement that the open subset be nonempty remain unresolved.
\subsection{Sources}
\begin{itemize}
\item Supplied local UnsolvedMath registry record: permanent ID \(701220\), exact source code \texttt{catalogue entry}, title \emph{Kourovka Notebook Problem 16.28}. Proposed by E. P. Vdovin in 2006.
\item \emph{The Kourovka Notebook (Unsolved Problems in Group Theory)}, 16th issue (2006), Problem 16.28; also cited in the supplied record through the 21st issue edition (2026), \texttt{arXiv:1401.0300v48}. The supplied collection description names E. I. Khukhro and V. D. Mazurov as editors.
\item Achyuth Jayadevan, \emph{Preprint of 12 September 2026}, \url{https://zenodo.org/records/22729229}. The supplied discussion reports a negative answer and describes this as an AI solution confirmed by E. P. Vdovin. This attribution and confirmation are supplied claims, not independently verified here.
\item M. Kinyon and J. Urban, \emph{Preprint of 17 September 2026}, \url{https://kourovkanotebookorg.wordpress.com/wp-content/uploads/2026/09/kourovka-mathematics.pdf}. The supplied discussion explicitly describes this as an unverified AI solution.
\item The registry labels the record \texttt{partially\_solved}. Its status review, dated exactly \texttt{2026-10-16}, reports that a subproblem is confirmed solved in the maintained 21st-issue source, while the background uses the label ``[Solved] No, not always.'' This assessment is supplied dated curation, not independently verified. The supplied excerpts do not identify which subproblem the negative answer resolves.
\end{itemize}
\end{document}
